\section{Problem Formulation}
\label{sec:problem}

\subsection{Modeling}
\label{sec:problem_modeling}

Our modeling objective is to capture the dynamics of a measurable LLM output attribute under prompt-based control, so that its future evolution can be predicted over multiple interaction steps.

At interaction step \(t\), a frozen large language model receives a task input, the accumulated dialogue history, and a textual control instruction, and produces the next response. The resulting input--output interaction is
\begin{equation}
o_{t+1}\sim
p_{\theta}(\cdot\mid x,\mathcal{H}_t,u_t),
\qquad
y_{t+1}=V(o_{t+1}),
\label{eq:interaction}
\end{equation}
Here \(x\) denotes the task input, \(p_{\theta}\) the frozen LLM, \(\mathcal{H}_t\) the dialogue history available at step \(t\), and \(u_t\) the textual instruction supplied by the controller. The response \(o_{t+1}\) is text, while \(V\) is a task-specific verifier that maps the completed response to a scalar attribute \(y_{t+1}\), such as sentence length or sentiment polarity. We track one such attribute at a time. A step corresponds to one complete prompt--response exchange.

The verifier output is the observable signal available to the controller. The internal activations that generate the response are inaccessible, so the measured attribute \(y_t\) is both the quantity available for system identification and the quantity in which the control objective is expressed.

A useful model must remain accurate beyond the next interaction, since control decisions depend on the consequences of an instruction over multiple future turns. We therefore evaluate the model through multi-step prediction error,
\begin{equation}
\ell_h
=
\bigl(y_{t+h}-\widehat y_{t+h\mid t}\bigr)^2,
\qquad h=1,\ldots,H.
\label{eq:modeling_objective}
\end{equation}
Here \(\widehat y_{t+h\mid t}\) denotes the prediction of the attribute at step \(t+h\) using information available at step \(t\), while \(y_{t+h}\) denotes the corresponding realized measurement. The prediction error should remain small throughout the horizon \(H\), rather than only for the one-step prediction \(h=1\). Section~\ref{sec:model} specifies how these forecasts are constructed and what information is held fixed during multi-step prediction.

\subsection{Control}
\label{sec:problem_control}

Our control objective is to regulate the measured output attribute toward a desired value through prompt-based commands within a finite interaction budget.

At each interaction step, the controller selects a scalar command \(c_t\), which determines the textual instruction \(u_t\) applied to the LLM. Over a finite horizon of \(H\) steps, the control objective is
\begin{equation}
J_t
=
\sum_{h=1}^{H}
\bigl(r-y_{t+h}\bigr)^2.
\label{eq:control_objective}
\end{equation}
Here \(r\in\mathbb{R}\) denotes the desired value of the tracked attribute, and \(e_t=r-y_t\) is the corresponding tracking error. The objective therefore seeks to keep the tracking error small over the interaction horizon, subject to the dynamics in Eq.~\eqref{eq:interaction}. In predictive control, the future measurements \(y_{t+h}\), which are unavailable at decision time, are replaced by their predictions \(\widehat y_{t+h\mid t}\) when selecting the control commands.

The controlled quantity is continuous, whereas the experiments ultimately report constraint satisfaction. We connect these two views through the task-specific feasible set \(\mathcal{Y}\): a response satisfies the constraint when \(y_t\in\mathcal{Y}\), and the reference is chosen such that \(r\in\mathcal{Y}\). Reducing the tracking error therefore provides a continuous control objective, while constraint satisfaction provides the corresponding task-level evaluation criterion. Driving \(y_t\) toward \(r\) is thus the mechanism through which the controller seeks to increase the rate of feasible responses.

The reference \(r\) and the command \(c_t\) have distinct roles. The reference specifies the desired output, whereas \(c_t\) is the action chosen by the controller and encoded as the instruction \(u_t\). A direct target-following controller uses \(c_t=r\). Predictive control instead allows \(c_t\neq r\) to compensate for the predicted response of the LLM. This requires a state representation that retains sufficient information from the interaction history to support accurate multi-step prediction.
